% Source: https://artofproblemsolving.com/community/c3h1342159_mathcounts
% Stable identity: aops-topic-1342159-q01
The second number $m$ satisfies $m+s(m)=44$, where $s$ is the digit sum. Since $s(m)\le 3+9=12$ for $m<44$, we have $32\le m\le43$. Checking: $m=40$ gives $40+4=44$, while $m=32,\dots,39$ give $37,39,41,43,45,47,49,51$ and $m=41,42,43$ give $46,48,50$. So $m=40$.

Now the first number $n$ satisfies $n+s(n)=40$, so $28\le n\le39$. Since $n+s(n)\equiv 2n\pmod 9$, we need $2n\equiv40\equiv4\pmod9$, so $n\equiv2\pmod9$, giving $n=29$ or $n=38$. And $29+11=40$ while $38+11=49$. The first number is $\boxed{29}$ (list: $29,40,44$).
